% ===========================================================================
% Section 7 — Conclusion and Future Work
% ===========================================================================

\section{Conclusion and Future Work}
\label{sec:conclusion}

% ---------------------------------------------------------------------------
\subsection{Summary of Contributions}
\label{sec:contributions}

This paper presented \textbf{SHARP-LLM} (Secure Healthcare-Aligned Risk
Prediction via Lightweight LLMs), a lightweight transformer-based framework
for automated CWE vulnerability classification in C/C++ source code with
healthcare-oriented risk prioritisation.  The framework makes five principal
contributions:

\begin{enumerate}
    \item \textbf{Template-aware data splitting.}
          We introduced a splitting strategy that partitions the NIST
          Juliet Test Suite by vulnerability template rather than by
          individual sample, eliminating structural data leakage between
          training and test sets.  Under this rigorous protocol, zero
          template overlap exists between the 82,736-sample training set
          (1,323~templates) and the 22,447-sample test set
          (366~templates), ensuring that reported metrics reflect genuine
          generalisation to unseen code patterns.

    \item \textbf{Multi-paradigm model evaluation.}
          We systematically compared three inference paradigms --- full
          encoder fine-tuning (CodeT5-Small, CodeT5-Base, CodeBERT-Base,
          GraphCodeBERT-Base), QLoRA-based parameter-efficient fine-tuning
          (CodeGemma-2B), and zero-shot generative inference
          (CodeGemma-2B, Gemma-4-E2B-IT) --- within a single unified
          pipeline.  To the best of our knowledge, this is the first study
          to conduct such a multi-paradigm comparison on 118-class CWE
          classification.

    \item \textbf{Consumer-grade GPU feasibility.}
          All experiments were conducted on a single NVIDIA RTX~2000~Ada
          laptop GPU with 8\,GB VRAM.  Fine-tuned encoder models train in
          under 53~minutes and achieve sub-30\,ms inference latency, while
          QLoRA enables fine-tuning of 2.5\,B-parameter models within the
          same memory budget through NF4 double quantisation, gradient
          checkpointing, and 8-bit paged optimisation.

    \item \textbf{Healthcare-oriented risk prioritisation pipeline.}
          Beyond CWE classification, the framework introduces a three-stage
          post-detection pipeline inspired by
          C3-VULMAP~\cite{santos2023c3vulmap}:
          (a)~\emph{privacy and security risk categorisation}, which maps
          each detected CWE to LINDDUN privacy threat types (linkability,
          identifiability, non-repudiation, detectability, information
          disclosure, unawareness, non-compliance) via a risk mapping
          function $\mathcal{R} = f(c)$;
          (b)~\emph{policy-relevant control domain mapping}, which
          associates risk categories with control areas such as
          confidentiality protection, access control, audit and
          accountability, and service availability via
          $\mathcal{C} = g(\mathcal{R})$; and
          (c)~\emph{weighted risk prioritisation}, which computes a
          composite score
          $P = \lambda_1 s + \lambda_2 E + \lambda_3 D + \lambda_4 S
          + \lambda_5 C'$
          combining model confidence~($s$), technical severity~($E$),
          data-sensitivity impact~($D$), service/safety impact~($S$),
          and control-domain importance~($C'$) to produce ordinal
          priority levels (Critical, High, Medium, Low).
          This pipeline transforms raw CWE predictions into actionable,
          healthcare-relevant risk signals for remediation planning.

    \item \textbf{Comprehensive evaluation and error analysis.}
          We reported nine metrics (accuracy, macro precision, recall, F1,
          MCC, FPR, model size, parameter count, and inference latency)
          across all models, analysed per-epoch convergence and overfitting
          dynamics, and performed per-class error analysis identifying
          systematic confusion patterns towards high-frequency buffer
          access CWEs.
\end{enumerate}

% ---------------------------------------------------------------------------
\subsection{Key Findings}
\label{sec:key_findings}

The experimental results (Section~\ref{sec:results}) yield several findings
with implications for the vulnerability detection community:

\begin{itemize}
    \item \textbf{Lightweight fine-tuning outperforms large-scale
          zero-shot models.}  CodeT5-Base (110\,M parameters) achieves a
          macro-F1 of 0.9605 on 118~CWE classes, surpassing the
          5.13\,B-parameter Gemma-4-E2B-IT (macro-F1 = 0.6988) by 0.2617
          in F1 while being $15{,}600\times$ faster at inference.  This
          demonstrates that task-specific fine-tuning of compact encoders
          remains the most effective strategy for structured vulnerability
          classification.

    \item \textbf{Encoder architecture matters more than parameter count.}
          CodeT5-Base (110\,M) outperforms both CodeBERT-Base and
          GraphCodeBERT-Base (each 125\,M) in macro-F1 (0.9605 vs.\ 0.9512
          and 0.9530) and exhibits more robust generalisation across
          training epochs.  The T5 encoder's bimodal pre-training objective
          produces superior representations for code vulnerability tasks
          compared to the RoBERTa-based architecture.

    \item \textbf{Early stopping is essential for RoBERTa-family models.}
          CodeBERT and GraphCodeBERT overfit at epoch~2, with macro-F1
          dropping by up to 5.78 percentage points.  A patience of
          1~epoch mitigates this degradation without sacrificing peak
          performance.

    \item \textbf{Instruction-tuning significantly improves zero-shot
          classification.}  Gemma-4-E2B-IT achieves $3\times$ the accuracy
          of CodeGemma-2B in zero-shot mode (78.0\% vs.\ 25.0\%),
          confirming that instruction alignment benefits structured
          classification tasks even without task-specific training data.

    \item \textbf{Template-aware evaluation prevents inflated metrics.}
          By ensuring no template overlap between train and test, our
          splitting strategy provides a more reliable estimate of model
          performance on genuinely novel code structures.  Despite this
          stricter protocol, all fine-tuned models still exceed 99.6\%
          accuracy, validating the approach.

    \item \textbf{Post-detection risk pipeline bridges code analysis
          and healthcare decision support.}
          The three-stage pipeline --- CWE-to-LINDDUN risk categorisation,
          policy-relevant control domain mapping, and weighted risk
          prioritisation --- converts raw classification outputs into
          prioritised, healthcare-relevant findings.  This structured
          interpretation layer enables remediation decisions aligned with
          patient-data confidentiality, service continuity, and regulatory
          compliance rather than treating all vulnerabilities as equally
          urgent.
\end{itemize}

% ---------------------------------------------------------------------------
\subsection{Broader Significance}
\label{sec:significance}

The results have practical implications for secure software development
pipelines.  The sub-30\,ms inference latency of fine-tuned encoder models
makes them viable for integration into CI/CD pipelines, IDE plugins, and
real-time code review tools.  The consumer-grade hardware requirement
(8\,GB VRAM) removes the need for expensive cloud GPU infrastructure,
lowering the barrier for adoption by small development teams, academic
institutions, and resource-constrained organisations.

The multi-class CWE classification capability is more actionable than binary
vulnerability detection: by identifying the \emph{specific} weakness type,
the system enables targeted remediation guidance aligned with
industry-standard taxonomies (MITRE CWE, OWASP).

Critically, SHARP-LLM goes beyond vulnerability detection by introducing a
healthcare-oriented post-detection pipeline.  The C3-VULMAP-inspired risk
categorisation stage~\cite{santos2023c3vulmap} maps each detected CWE to
LINDDUN privacy threat types, enabling the framework to distinguish between
weaknesses that affect patient-data confidentiality, traceability, service
continuity, or compliance obligations.  The subsequent policy-relevant
control domain mapping and weighted risk prioritisation stages translate
these risk signals into ordinal priority levels (Critical / High / Medium /
Low), providing security analysts with a structured basis for triage that
reflects the unique concerns of healthcare software environments.  This
three-stage pipeline positions SHARP-LLM not merely as a vulnerability
detector but as a healthcare-aligned decision support tool.

% ---------------------------------------------------------------------------
\subsection{Limitations}
\label{sec:limitations}

Despite the strong results, the current work has several limitations that
must be acknowledged:

\begin{enumerate}
    \item \textbf{Synthetic dataset.}
          All experiments use the NIST Juliet Test Suite, which contains
          template-generated synthetic code.  While the template-aware split
          mitigates trivial pattern memorisation, the code lacks the
          diversity, complexity, and contextual richness of real-world
          software projects.  Models trained exclusively on Juliet may not
          generalise to production codebases with multi-file dependencies,
          third-party libraries, and non-standard coding styles.

    \item \textbf{Class imbalance and single-sample categories.}
          The dataset exhibits a long-tailed distribution with an imbalance
          ratio exceeding 1,500:1.  Thirty-one CWE categories contain only
          a single template and are therefore absent from the test set,
          precluding evaluation on those classes.  Three categories with
          single test samples yield F1~=~0.00, but this reflects
          insufficient evaluation data rather than model failure.

    \item \textbf{Sequence length constraint.}
          All models truncate input at 256~tokens.  While sufficient for the
          Juliet samples (which are typically short, self-contained
          functions), this limit may cause information loss on longer
          real-world source files where the vulnerability context spans
          hundreds of lines.

    \item \textbf{Zero-shot evaluation on limited samples.}
          Due to prohibitive inference costs (up to 396\,s per sample for
          Gemma-4-E2B-IT), zero-shot models were evaluated on sampled
          subsets of 50--500 samples.  Larger-scale evaluation may yield
          different aggregate statistics and is necessary for definitive
          conclusions about zero-shot performance.

    \item \textbf{Single dataset and language.}
          The framework has been evaluated only on C/C++ code from the
          Juliet Test Suite.  Vulnerability patterns in other languages
          (Java, Python, JavaScript) and from other data sources may
          present different challenges for the models.

    \item \textbf{No cross-project evaluation.}
          All training and evaluation occur within the Juliet ecosystem.
          Cross-project generalisation --- where the model is trained on one
          codebase and tested on a different one --- remains untested.

    \item \textbf{Risk prioritisation weights are not empirically
          calibrated.}
          The weighted risk prioritisation score
          ($P = \lambda_1 s + \lambda_2 E + \lambda_3 D + \lambda_4 S
          + \lambda_5 C'$) relies on configurable coefficients that must
          be tuned to the deployment environment.  The current work
          presents the formulation but does not include an empirical
          calibration study with healthcare domain experts to determine
          optimal weight settings.

    \item \textbf{C3-VULMAP coverage.}
          The CWE-to-LINDDUN risk mapping derived from
          C3-VULMAP~\cite{santos2023c3vulmap} covers a defined subset of
          healthcare-relevant CWEs.  Detected CWE categories that fall
          outside this mapping currently receive a default risk
          categorisation, which may under- or over-estimate their
          healthcare impact.
\end{enumerate}

% ---------------------------------------------------------------------------
\subsection{Future Work}
\label{sec:future_work}

We identify the following concrete directions for extending this research:

\subsubsection{Data-Centric Improvements}

\begin{enumerate}
    \item \textbf{Real-world dataset augmentation.}
          Incorporating samples from established real-world vulnerability
          datasets such as Big-Vul~\cite{fan2020bigvul},
          Devign~\cite{zhou2019devign}, ReVeal~\cite{chakraborty2021reveal},
          and D2A~\cite{zheng2021d2a} would expose models to genuine
          vulnerability patterns beyond synthetic templates.  A mixed
          training regime combining Juliet and real-world data could
          improve generalisation while retaining the controlled evaluation
          benefits of the template-aware split.

    \item \textbf{Semantic-preserving code augmentation.}
          Applying transformations such as variable renaming, dead code
          insertion, control flow restructuring, and comment
          removal would force models to learn deeper vulnerability
          semantics rather than surface-level syntactic patterns.  This
          is particularly relevant for mitigating the CWE-126 confusion
          bias identified in the error analysis
          (Section~\ref{sec:performance}).

    \item \textbf{LLM-generated synthetic diversity.}
          Large language models (e.g., GPT-4, CodeLlama) could generate
          diverse vulnerable code variants for each CWE category, providing
          additional training samples with varied coding styles and contexts
          beyond the Juliet templates.
\end{enumerate}

\subsubsection{Architecture and Training Enhancements}

\begin{enumerate}
    \setcounter{enumi}{3}
    \item \textbf{Contrastive learning on good/bad pairs.}
          The Juliet Test Suite contains paired vulnerable and patched
          samples for each template.  Training with a contrastive objective
          on these pairs would teach the model to discriminate
          \emph{what makes code vulnerable} rather than simply classifying
          surface patterns, potentially reducing inter-class confusion.

    \item \textbf{Multi-task learning.}
          Adding a secondary binary classification head (vulnerable vs.\
          safe) alongside the 118-class CWE head would provide additional
          regularisation and could improve feature learning for
          underrepresented vulnerability categories.

    \item \textbf{Graph-enhanced representations.}
          Incorporating Abstract Syntax Tree (AST), Control Flow Graph
          (CFG), or Program Dependence Graph (PDG) information alongside
          token sequences could capture structural vulnerability patterns
          that purely sequence-based models miss.  Hybrid architectures
          combining GNN encoders with transformer embeddings represent a
          promising direction.

    \item \textbf{Ensemble with confidence calibration.}
          Combining predictions from multiple fine-tuned models (CodeT5-Base,
          CodeBERT, GraphCodeBERT) via temperature-scaled ensembling could
          improve accuracy on ambiguous samples.  Model disagreement could
          serve as an uncertainty signal for flagging samples that require
          human review.
\end{enumerate}

\subsubsection{Deployment and Scalability}

\begin{enumerate}
    \setcounter{enumi}{7}
    \item \textbf{Cross-language and cross-project evaluation.}
          Extending the framework to Java, Python, and JavaScript
          vulnerability datasets would test the transferability of the
          approach.  Cross-project evaluation, where models train on one
          codebase and are tested on an entirely different project, would
          provide the most rigorous assessment of real-world applicability.

    \item \textbf{Integration with static analysis tools.}
          A hybrid approach combining neural predictions with traditional
          static analysis features (taint paths, buffer sizes, pointer
          aliases from tools such as Cppcheck, Semgrep, or CodeQL) could
          provide the model with domain knowledge that is difficult to learn
          from token sequences alone, particularly for complex
          vulnerability patterns.

    \item \textbf{Longer context and function-level slicing.}
          Increasing the maximum sequence length beyond 256~tokens or
          employing vulnerability-relevant program slicing (extracting only
          the code regions involved in data/control flow to the vulnerable
          point) would reduce noise and enable analysis of larger,
          multi-function source files encountered in production codebases.

    \item \textbf{QLoRA fine-tuning completion and scaling.}
          Completing the QLoRA fine-tuning of CodeGemma-2B and extending the
          approach to larger instruction-tuned models (e.g., Gemma-4-E2B-IT)
          would quantify the performance gains of parameter-efficient
          fine-tuning over zero-shot inference for decoder-based models,
          bridging the gap between the encoder and zero-shot paradigms
          evaluated in this study.
\end{enumerate}

\subsubsection{Healthcare Risk Pipeline Refinement}

\begin{enumerate}
    \setcounter{enumi}{11}
    \item \textbf{Empirical calibration of prioritisation weights.}
          Conducting a user study with healthcare security analysts to
          empirically calibrate the risk prioritisation coefficients
          ($\lambda_1$--$\lambda_5$) would ground the scoring model in
          real-world triage practices and improve its alignment with
          clinical and operational priorities.

    \item \textbf{Extended CWE-to-LINDDUN mapping.}
          Expanding the C3-VULMAP-inspired mapping to cover additional CWE
          categories --- particularly those related to cryptographic
          failures, authentication weaknesses, and supply-chain risks ---
          would increase the coverage and accuracy of the risk
          categorisation stage for diverse healthcare software ecosystems.

    \item \textbf{Regulatory framework integration.}
          Linking the policy-relevant control domains to specific
          regulatory requirements (e.g., HIPAA Security Rule safeguards,
          GDPR data protection principles, IEC~62443 for medical device
          cybersecurity) would enable the framework to produce
          compliance-oriented reports, further strengthening its utility
          as a healthcare decision support tool.
\end{enumerate}

\subsection*{}
In summary, the SHARP-LLM framework demonstrates that lightweight,
fine-tuned transformer models achieve state-of-the-art multi-class CWE
classification with macro-F1 exceeding 0.96 on 118~vulnerability
categories, while training in under one hour on a single consumer-grade
GPU.  By augmenting vulnerability detection with a healthcare-oriented
risk prioritisation pipeline --- mapping detected CWEs to LINDDUN privacy
threat types, policy-relevant control domains, and ordinal priority
levels --- the framework bridges the gap between source code analysis and
healthcare decision support.  These results establish a strong foundation
for practical, cost-effective automated vulnerability detection and risk
triage that can be integrated into modern secure software development
workflows for healthcare and other safety-critical domains.
